% ════════════════════════════════════════════════════════════════
%  DNB BLANC N°2 — Collège Gaston Chaissac, Pouzauges
%  Session du jeudi 9 avril 2026
%  Compilé avec : pdflatex dnb_blanc_2_2026.tex
% ════════════════════════════════════════════════════════════════

\documentclass[a4paper,11pt]{article}

%% ─── Encodage et langue ──────────────────────────────────────
\usepackage[utf8]{inputenc}
\usepackage[T1]{fontenc}
\usepackage[french]{babel}
%% Formatage des nombres : espace fine comme séparateur des milliers
\newcommand{\np}[1]{\numprint{#1}}
\usepackage[autolanguage]{numprint}

%% ─── Police professionnelle ──────────────────────────────────
\usepackage{fourier}

%% ─── Mise en page ────────────────────────────────────────────
\usepackage[top=2cm, bottom=2cm, left=2cm, right=2cm]{geometry}
\usepackage{fancyhdr}
\pagestyle{fancy}
\fancyhf{}
\cfoot{\thepage}
\renewcommand{\headrulewidth}{0pt}
\setlength{\parindent}{0pt}

%% ─── Mathématiques ───────────────────────────────────────────
\usepackage{amsmath, amssymb}

%% ─── Listes ──────────────────────────────────────────────────
\usepackage{enumitem}

%% ─── Tableaux ────────────────────────────────────────────────
\usepackage{tabularx, array, multirow}

%% ─── Graphiques ──────────────────────────────────────────────
\usepackage{tikz}
\usetikzlibrary{arrows.meta, calc, positioning}
\usepackage{pgfplots}
\pgfplotsset{compat=1.18}

%% ─── Boîtes ──────────────────────────────────────────────────
\usepackage{tcolorbox}
\tcbuselibrary{skins}

%% ─── Scratch ─────────────────────────────────────────────────
\usepackage{scratch3}

%% ─── Divers ──────────────────────────────────────────────────
\usepackage{xcolor}

%% ─── Commande exercice ───────────────────────────────────────
\newcommand{\exercice}[2]{%
  \newpage
  \noindent{\Large\bfseries Exercice #1\hspace{0.3em}{\normalfont\normalsize(#2 points)}}%
  \par\vspace{0.6cm}%
}

% ══════════════════════════════════════════════════════════════
\begin{document}
\thispagestyle{fancy}

% ──────────────────────────────────────────────────────────────
%  PAGE DE GARDE
% ──────────────────────────────────────────────────────────────

\begin{flushleft}
{\fontsize{12}{14}\selectfont\textit{Collège Gaston CHAISSAC}}\\[2pt]
{\fontsize{12}{14}\selectfont\textit{POUZAUGES}}
\end{flushleft}

\vspace{1.5cm}

\begin{center}
  {\LARGE\bfseries DIPLÔME NATIONAL DU BREVET}\\[0.6cm]
  {\large SESSION DU JEUDI 9 AVRIL 2026}\\[0.3cm]
  {\large DURÉE : 2\,H \;--\; 100 POINTS}
\end{center}

\vspace{1cm}

% Encadré central gris
\begin{center}
\begin{tcolorbox}[
  width=12cm,
  colback=gray!25,
  colframe=gray!25,
  arc=4pt,
  boxrule=0pt,
  halign=center,
  top=10pt, bottom=10pt
]
  {\normalsize Épreuve de :}\\[0.5cm]
  {\Large\bfseries MATHÉMATIQUES}\\[0.4cm]
  {\large\itshape SERIE DNB BLANC N°2}\\[0.5cm]
  % Badge blanc
  \begin{tcolorbox}[
    width=7cm,
    colback=white,
    colframe=black,
    arc=3pt,
    boxrule=0.5pt,
    halign=center,
    top=4pt, bottom=4pt
  ]
    {\normalsize\bfseries BREVET BLANC}
  \end{tcolorbox}
\end{tcolorbox}
\end{center}

\vspace{1cm}

\begin{center}
  L'utilisation de la calculatrice est autorisée (circ.\ 99-186 du 16 novembre 1999)\\[0.3cm]
  Le sujet est composé de six exercices indépendants.\\[0.15cm]
  Le candidat peut les traiter dans l'ordre qui lui convient.
\end{center}

\vspace{0.8cm}

% Tableau barème
\begin{center}
\renewcommand{\arraystretch}{1.4}
\begin{tabular}{|p{5cm}|p{3cm}|}
  \hline
  Exercice 1 & 6 points  \\ \hline
  Exercice 2 & 4 points  \\ \hline
  Exercice 3 & 5 points  \\ \hline
  Exercice 4 & 7 points  \\ \hline
  Exercice 5 & 8 points  \\ \hline
  Exercice 6 & 5 points  \\ \hline
  \multicolumn{2}{|c|}{\small Barème indicatif} \\ \hline
\end{tabular}
\end{center}

\vfill

% Encadré instructions
\begin{center}
\fbox{\parbox{0.9\textwidth}{%
  \textbf{Toutes les réponses doivent être justifiées,}\quad sauf si une indication contraire est donnée.\\[0.2cm]
  Pour chaque question, si le travail n'est pas terminé, laisser tout de même une trace de la recherche.\\[0.1cm]
  Elle sera prise en compte dans la notation.%
}}
\end{center}

\vspace{0.5cm}

% ──────────────────────────────────────────────────────────────
%  EXERCICE 1 — Automatismes
% ──────────────────────────────────────────────────────────────
\exercice{1}{6}

\begin{enumerate}[label=\textbf{\arabic*)}, leftmargin=1.2cm, itemsep=0.5cm]

\item ABCD est un carré et CDE est un triangle équilatéral (E est à l'intérieur du carré ABCD).\\
BCF est un triangle isocèle en F (F est à l'extérieur du carré ABCD).\\
Représenter cette configuration par un schéma à main levée et ajouter les codages nécessaires.

\item Calculer le périmètre d'un cercle de rayon 3~cm.

\item Calculer : $77{,}6 \div 10$

\item Le diagramme suivant représente les ventes d'un magasin sur une semaine :

\begin{center}
\begin{tikzpicture}
\begin{axis}[
  ybar,
  bar width=1.2cm,
  symbolic x coords={Lundi,Mardi,Mercredi,Jeudi,Vendredi},
  xtick=data,
  nodes near coords,
  nodes near coords style={font=\bfseries\small},
  ymin=0, ymax=58,
  ytick={0,10,...,50},
  width=11cm, height=6.5cm,
  enlarge x limits=0.15,
  axis lines*=left,
  ymajorgrids=false,
  axis line style={thick},
  tick style={thick},
  every tick label/.style={font=\small},
  fill=green!50!black,
  draw=green!70!black,
]
\addplot coordinates {(Lundi,33) (Mardi,49) (Mercredi,48) (Jeudi,30) (Vendredi,32)};
\end{axis}
\end{tikzpicture}
\end{center}

Quel est le total des ventes sur la semaine ?

\item Une ville compte \np{49298} habitants pour une superficie de 302~km$^2$. Quelle est sa densité de population ?

\item Écrire un script Scratch qui dessine un polygone régulier à 5 côtés de 147 pixels de côté.

\end{enumerate}

% ──────────────────────────────────────────────────────────────
%  EXERCICE 2 — Circuits d'entraînement
% ──────────────────────────────────────────────────────────────
\exercice{2}{4}

Un entraîneur de sport prépare deux circuits d'entraînement contenant plusieurs exercices de cardio et de renforcement musculaire :

\begin{itemize}[label=\textbullet, itemsep=2pt]
  \item un circuit commence à l'exercice 1 et se termine en revenant à l'exercice 1 ;
  \item le circuit 1 contient cinq exercices. Chaque exercice dure 40 secondes et doit être suivi de 16 secondes de repos permettant de se rendre à l'exercice suivant ;
  \item le circuit 2 contient dix exercices. Chaque exercice dure 30 secondes et doit être suivi de 5 secondes de repos permettant de se rendre à l'exercice suivant.
\end{itemize}

\begin{center}
\begin{tikzpicture}[line cap=round,baseline={(C)}]
  \draw[<-, >=stealth] (0,2.9)--(0,3.5) node[above]{Départ / Arrivée};
  \foreach \ex in {1,...,5}{
    \draw[line width=1.5pt] (162-\ex*72:2) circle (0.2);
    \node at (162-\ex*72:2.7) {Ex~\ex};
    \draw (162-\ex*72:2)--(90-\ex*72:2);
  }
  \node (C) at(0,0) {Circuit 1};
\end{tikzpicture}
\hfill
\begin{tikzpicture}[line cap=round,baseline={(C)}]
  \draw[<-, >=stealth] (0,2.9)--(0,3.5) node[above]{Départ / Arrivée};
  \foreach \ex in {1,...,10}{
    \draw[line width=1.5pt] (126-\ex*36:2) circle (0.2);
    \node at (126-\ex*36:2.7) {Ex~\ex};
    \draw (126-\ex*36:2)--(90-\ex*36:2);
  }
  \node (C) at(0,0) {Circuit 2};
\end{tikzpicture}
\end{center}

\begin{enumerate}[itemsep=4pt]
  \item Montrer que le circuit 1 s'effectue en 280 secondes et que le circuit 2 s'effectue en 350 secondes.
  \item Donner la décomposition en produit de facteurs premiers de $280$ et de $350$.
  \item Une séance d'entraînement est constituée de plusieurs tours du même circuit.

  Au coup de sifflet de l'entraîneur, Camille commence une séance d'entraînement sur le circuit~1 et Dominique sur le circuit~2.
  \begin{enumerate}[label=\alph*., itemsep=4pt]
    \item Expliquer pourquoi, lorsque \np{2800} secondes se sont écoulées à partir du coup de sifflet, Camille se trouve de nouveau au départ du circuit~1.

    Préciser où se trouve Dominique sur le circuit~2 lorsque \np{2800} secondes se sont écoulées.

    \item Après le coup de sifflet, combien de temps faut-il à Camille et Dominique pour se retrouver en même temps pour la première fois au départ de leur circuit ? Exprimer cette durée en minute et seconde.
  \end{enumerate}
\end{enumerate}

% ──────────────────────────────────────────────────────────────
%  EXERCICE 3 — Programme de calcul
% ──────────────────────────────────────────────────────────────
\exercice{3}{5}

On considère le programme de calcul ci-dessous :

\begin{center}
\begin{tabular}{|l|}
  \hline
  $\bullet~$Choisir un nombre \\
  $\bullet~$Mettre ce nombre au carré \\
  $\bullet~$Soustraire le triple du nombre de départ \\
  $\bullet~$Soustraire 4 \\
  \hline
\end{tabular}
\end{center}

\smallskip

\begin{enumerate}[itemsep=4pt]
  \item Montrer que si on choisit 5 comme nombre de départ, le résultat du programme est 6.
  \item On choisit $x$ comme nombre de départ.

  Exprimer le résultat du programme en fonction de $x$.
  \item Vérifier que l'on peut écrire ce résultat sous la forme $(x + 1)(x - 4)$.
  \item Déterminer les nombres à choisir au départ pour que le résultat du programme soit 0.
  \item Juliette a écrit le programme ci-dessous :

\begin{center}
\begin{scratch}[num blocks]
\blockinit{quand \greenflag est cliqué}
\blockmove{demander \ovalnum{Choisir un nombre} et attendre}
\blockvariable{mettre \selectmenu{x} à \ovalmove{réponse}}
\blockvariable{mettre \selectmenu{y} à \ovaloperator{\ovalnum{\ldots}\,*\,\ovalnum{\ldots}}}
\blockvariable{mettre \selectmenu{z} à \ovaloperator{\ovalnum{3}\,*\,\ovalnum{x}}}
\blockvariable{mettre \selectmenu{Résultat} à \ovaloperator{\ovalnum{\ldots}\,-\,\ovalnum{\ldots}\,- 4}}
\blocklook{dire \ovalnum{Résultat} pendant \ovalnum{5} secondes}
\end{scratch}
\end{center}
\end{enumerate}

\medskip

Recopier et compléter sur la copie les lignes 4 et 6 du programme afin que celui-ci corresponde au programme de calcul encadré.

% ──────────────────────────────────────────────────────────────
%  EXERCICE 4 — Jardin botanique
% ──────────────────────────────────────────────────────────────
\exercice{4}{7}

Le jardin botanique d'une ville peut être représenté par le quadrilatère ABCD ci-dessous.

\medskip

\begin{minipage}{0.52\linewidth}
\begin{tikzpicture}[scale=0.85, thick]
  % Coordonnées
  \coordinate (A) at (0,4.33);
  \coordinate (B) at (2.5,4.33);
  \coordinate (C) at (7.5,0);
  \coordinate (D) at (0,0);
  \coordinate (E) at (1.875,3.248);

  % Remplissage gris
  \fill[gray!25] (A) -- (B) -- (E) -- cycle;
  \fill[gray!25] (D) -- (E) -- (C) -- cycle;

  % Quadrilatère
  \draw (A) -- (B) -- (C) -- (D) -- cycle;

  % Diagonales (pointillés)
  \draw[densely dotted, line width=1.2pt] (A) -- (C);
  \draw[densely dotted, line width=1.2pt] (B) -- (D);

  % Angles droits
  \draw (0.25,4.33) -- (0.25,4.08) -- (0,4.08);
  \draw (0.25,0) -- (0.25,0.25) -- (0,0.25);
  % Angle droit en E (rotation ~18.4°)
  \begin{scope}[shift={(E)}, rotate=-71.6]
    \draw (0.2,0) -- (0.2,0.2) -- (0,0.2);
  \end{scope}

  % Labels
  \node[above left]  at (A) {A};
  \node[above right] at (B) {B};
  \node[below right] at (C) {C};
  \node[below left]  at (D) {D};
  \node[left=4pt]    at (E) {E};
\end{tikzpicture}
\end{minipage}%
\hfill
\begin{minipage}{0.45\linewidth}
On sait que :

\smallskip

$\bullet\;$ AB $= 500$~m, BE $= 250$~m et DE $= 750$~m ;

\smallskip

$\bullet\;$ les segments [AC] et [BD] se coupent au point E.

\smallskip

\emph{La figure ci-contre n'est pas à l'échelle.}
\end{minipage}

\bigskip

\begin{enumerate}[itemsep=4pt]
  \item Quelle est la longueur du segment [DB]\,?
  \item En raisonnant dans le triangle rectangle ABD, montrer que la longueur du segment [AD], arrondie au mètre, est égale à environ $866$~m.
  \item
  \begin{enumerate}[label=\alph*., itemsep=3pt]
    \item Calculer le sinus de l'angle $\widehat{\text{EAB}}$.
    \item En déduire la mesure en degrés de l'angle $\widehat{\text{EAB}}$.
  \end{enumerate}
  \item
  \begin{enumerate}[label=\alph*., itemsep=3pt]
    \item Montrer que les droites (AB) et (DC) sont parallèles.
    \item Montrer que la longueur du segment [CD] est égale à \np{1500}~m.
  \end{enumerate}
  \item Un piéton fait le tour du jardin botanique en marchant à la vitesse moyenne de 1,1~m/s.

  Il lit sur son plan que la longueur du segment [BC] est environ égale à \np{1323}~m.

  Le temps mis par le piéton pour faire le tour du jardin botanique est-il inférieur à une heure\,?
\end{enumerate}

% ──────────────────────────────────────────────────────────────
%  EXERCICE 5 — Fonctions
% ──────────────────────────────────────────────────────────────
\exercice{5}{8}

\begin{minipage}[t]{7.2cm}
\begin{enumerate}
\item On considère le programme A défini par le schéma ci-contre :
  \begin{enumerate}[label=\alph*., itemsep=3pt]
    \item Vérifier que le résultat est $60$ si le nombre choisi au départ est $-8$.
    \item On appelle $x$ le nombre de départ et on admet que le résultat obtenu avec le programme de calcul est donné par l'expression :

    $(x + 3)(x - 4)$.

    Résoudre $(x + 3)(x - 4) = 0$.

    En déduire quels nombres de départ il faut choisir pour obtenir 0 comme résultat.
  \end{enumerate}
\end{enumerate}
\end{minipage}\hfill
\begin{minipage}[t]{7.5cm}
\vspace{0pt}
\begin{center}
\begin{tikzpicture}[
  block/.style={draw, thick, rectangle, minimum width=1.4cm, minimum height=0.6cm},
  arr/.style={->, >=Stealth, thick},
  font=\small,
  scale=0.85
]
  \node at (0,7.1) {Choisir un nombre};
  \draw[arr] (0,6.8) -- (0,6.2);
  \node[block] (top) at (0,5.8) {};

  \node at (-2.3,4.8) {Ajouter 3};
  \node at (2.7,4.8) {Soustraire 4};

  \draw[arr] (0,5.4) -- (-2.9,3.9);
  \draw[arr] (0,5.4) -- (2.9,3.9);

  \node[block] (L) at (-2.9,3.5) {};
  \node[block] (R) at (2.9,3.5) {};

  \draw (L.south) -- (0,1.4);
  \draw (R.south) -- (0,1.4);

  \node[align=center] at (0,2.45) {Multiplier les deux\\résultats};

  \draw[arr] (0,1.4) -- (0,0.8);
  \node[block] (bot) at (0,0.4) {};
\end{tikzpicture}
\end{center}
\end{minipage}

\medskip

\begin{enumerate}[resume]
\item On rappelle que $x$ désigne le nombre de départ du programme de calcul et que le résultat obtenu avec le programme de calcul est donné par l'expression : $(x + 3)(x - 4)$.

On appelle $f$ la fonction qui, à $x$, associe le résultat du programme de calcul.

La représentation graphique $\mathcal{C}_f$ de la fonction $f$ est donnée ci-après.
  \begin{enumerate}[label=\alph*., itemsep=3pt]
    \item Montrer que $f(x) = x^2 - x - 12$.
    \item Calculer $f\!\left(\dfrac{1}{2}\right)$.
    \item Déterminer graphiquement les antécédents de $-6$ par la fonction $f$.

    On pourra éventuellement laisser les traits de construction sur \textbf{le graphique}.
  \end{enumerate}
\end{enumerate}

\begin{center}
\begin{tikzpicture}
\begin{axis}[
  axis lines=middle,
  xmin=-5, xmax=6.2,
  ymin=-12.9, ymax=12,
  xtick={-4,-3,...,5},
  ytick={-12,-11,...,11},
  minor tick num=0,
  grid=major,
  grid style={gray!30, thin},
  width=12cm,
  height=14cm,
  axis line style={thick, ->, >=Stealth},
  every tick label/.style={font=\scriptsize},
  xlabel={}, ylabel={},
  clip=true,
]
\addplot[blue, thick, smooth, domain=-4.6:5.6, samples=300] {x^2 - x - 12};
\node[blue, left, font=\small] at (axis cs:-3.8,10) {$\mathcal{C}_f$};
\end{axis}
\end{tikzpicture}
\end{center}

\begin{minipage}[t]{10cm}
\begin{enumerate}[resume]
\item On considère la fonction $g$ définie par $g(x) = 3x - 7$.

On a utilisé un tableur pour réaliser un tableau de valeurs de cette fonction.
  \begin{enumerate}[label=\alph*., itemsep=3pt]
    \item Quelle formule a-t-on écrite dans la cellule B2 avant de l'étirer vers le bas\,?
    \item Tracer la représentation graphique de la fonction $g$ dans le repère sur \textbf{le graphique précédent}.
    \item Déterminer graphiquement les nombres qui ont la même image par les fonctions $f$ et $g$. On pourra laisser apparents les traits de construction sur \textbf{le graphique}.
  \end{enumerate}
\end{enumerate}
\end{minipage}\hfill
\begin{minipage}[t]{4.5cm}
\vspace{0pt}
\renewcommand{\arraystretch}{1.15}
\begin{center}
\begin{tabular}{|c|c|c|}
  \hline
     & A    & B       \\ \hline
  1  & $x$  & $g(x)$  \\ \hline
  2  & $-5$ & $-22$   \\ \hline
  3  & $-4$ & $-19$   \\ \hline
  4  & $-3$ & $-16$   \\ \hline
  5  & $-2$ & $-13$   \\ \hline
  6  & $-1$ & $-10$   \\ \hline
  7  & $0$  & $-7$    \\ \hline
  8  & $1$  & $-4$    \\ \hline
  9  & $2$  & $-1$    \\ \hline
  10 & $3$  & $2$     \\ \hline
  11 & $4$  & $5$     \\ \hline
  12 & $5$  & $8$     \\ \hline
  13 & $6$  & $11$    \\ \hline
\end{tabular}
\end{center}
\end{minipage}

% ──────────────────────────────────────────────────────────────
%  EXERCICE 6 — Lunettes de soleil
% ──────────────────────────────────────────────────────────────
\exercice{6}{5}

Un opticien vend différents modèles de lunettes de soleil.

Il reporte dans le tableur ci-dessous des informations sur cinq modèles vendus pendant l'année 2022.

\medskip

\begin{center}
\renewcommand{\arraystretch}{1.4}
\begin{tabularx}{\linewidth}{|c|>{\centering\arraybackslash}m{3.5cm}|*{6}{>{\centering\arraybackslash}X|}}
  \hline
     & \textsf{A} & \textsf{B} & \textsf{C} & \textsf{D} & \textsf{E} & \textsf{F} & \textsf{G} \\ \hline
  \textsf{1} & \textbf{Lunettes de soleil} & \textbf{Modèle 1} & \textbf{Modèle 2} & \textbf{Modèle 3} & \textbf{Modèle 4} & \textbf{Modèle 5} & \textbf{Total} \\ \hline
  \textsf{2} & \textbf{Nombre de paires de lunettes vendues} & \np{1200} & 950 & 875 & 250 & 300 &  \\ \hline
  \textsf{3} & \textbf{Prix à l'unité en euro} & 75 & 100 & 110 & 140 & 160 &  \\ \hline
\end{tabularx}
\end{center}

\medskip

\begin{enumerate}[itemsep=4pt]
  \item Montrer que l'étendue des prix de ces paires de lunettes de soleil est de 85 euros.
  \item
  \begin{enumerate}[label=\alph*., itemsep=3pt]
    \item Quelle formule doit-on saisir dans la cellule G2 pour calculer le nombre total de paires de lunettes de soleil vendues en 2022\,?
    \item Calculer le nombre total de paires de lunettes de soleil vendues en 2022.
  \end{enumerate}
  \item
  \begin{enumerate}[label=\alph*., itemsep=3pt]
    \item Calculer le montant total, en euros, des ventes des paires de lunettes de soleil en 2022.
    \item Calculer le prix moyen d'une paire de lunettes de soleil vendue en 2022, arrondi au centime près.
  \end{enumerate}
\end{enumerate}

\end{document}
